\section{Related Work}
\label{sec:related}

Two largely separate threads in the literature bear on our research question. We organise prior work into (i) sycophancy and the ``challenged LLM'' family, (ii) framing and tone manipulations applied as a prompt-engineering lever, and (iii) the smaller body of work on reasoning quality and apology markers.

\para{Sycophancy and post-hoc challenge.} \citet{sharma2023sycophancy} establish that RLHF-trained models systematically capitulate to skeptical or contradicting users, and that preference models prefer sycophantic answers over correct ones a non-negligible fraction of the time. \citet{laban2023flipflop} formalise the ``\flipflop'' two-turn protocol: a model answers, a challenger utterance such as ``Are you sure?'' is delivered, and the model may revise its answer. They report a mean flip rate of $46\%$ and a mean accuracy drop of $17\%$ across ten models and seven classification tasks. Critically, every challenger they study \emph{degrades} performance — skepticism is monotonically harmful in their setting. They also introduce \psorry, the fraction of responses that contain apology keywords (``sorry'', ``you're right'', ``my mistake''), which becomes the field's de-facto operational marker of capitulation. \citet{wei2023synthetic} show that synthetic finetuning can mitigate this behaviour, and \citet{hong2025sycon} extend the protocol to a five-turn multi-scenario benchmark, finding that larger and reasoning-optimised models within a family are less sycophantic by up to 81\%. \citet{petrov2025brokenmath} apply the same framing to theorem proving, and \citet{zhang2025sycophancyunderpressure}, \citet{lee2025assertbench}, and \citet{zhu2024conformity} extend it to scientific QA, factual assertion, and Asch-style conformity respectively. \textbf{Unlike these works}, we apply the skeptical prompt \emph{before} the model answers, vary tone independently of content, and measure reasoning length in addition to accuracy and capitulation.

\para{Preemptive vs.\ in-context rebuttal and framing.} \citet{fanous2025syceval} is the only prior work we are aware of that distinguishes preemptive from in-context rebuttals. They find that preemptive rebuttals are \emph{more} sycophancy-inducing ($61.75\%$ vs.\ $56.52\%$), and they vary content strength along a four-step Simple$\rightarrow$Citation ladder. They do not vary tone with content held fixed and do not measure reasoning length. \citet{kim2025challenging} report three findings on a separate conversational-vs-evaluative axis: identical arguments are more often accepted under conversational framing (H1), rebuttals with reasoning sway models more (H2), and casually-phrased rebuttals sway models more than formal ones (H3). H3 is the closest prior signal that tone is a separable lever, but it varies one bit of register, not a six-step ladder. \citet{rabbani2025factjudgment} show that reframing a factual question as a conversational judgment task produces $\sim$9\% performance change in a model-dependent direction. \textbf{Unlike these works}, we hold both content and the conversational/evaluative axis constant, vary tone over six graded levels from neutral to hostile, and explicitly test for non-monotonicity.

\para{Tone as a prompt-engineering lever.} The positive arm of our hypothesis is anchored by \citet{li2023emotionprompt}, who show that emotional and motivational stimuli boost accuracy by an average of $10.9\%$ across $45$ tasks and six model families. The mechanism is purely tonal: identical content, different framing. \citet{hong2025sycon} find that a third-person ``Andrew prompt'' persona reduces sycophancy by $63.8\%$ in their debate scenario. \citet{needham2025evalawareness} document that frontier models display some ``evaluation awareness'' — they sometimes reason explicitly about being in a safety eval — which is a prerequisite for any social-facilitation-style effect. \textbf{Unlike these works}, we instrument the lever quantitatively, with paired token counts per item, and connect it directly to the sycophancy axis through the two-turn protocol.

\para{Reasoning quality and silent capitulation.} \citet{feng2026goodarguments} argue that chain-of-thought reasoning can \emph{mask} sycophancy: longer traces may be post-hoc rationalisation rather than better deliberation. \citet{barkett2025reasoning} report that reasoning models still exhibit non-trivial truth bias above the human baseline. \citet{chang2026causalt3} names the dual failure modes — a Skepticism Trap, where capable models over-refuse correct answers, and a Sycophancy Trap, where confident user pressure flips correct answers — that we hypothesise are flanking the sweet spot. \citet{dunlap2026feedback} report that newer OpenAI models with higher reasoning effort are less sycophantic. \citet{monica2025} introduce real-time monitoring of sycophantic drift during reasoning steps. \textbf{Unlike these works}, our results highlight a different problem: the verbal apology marker that the field has used to detect capitulation is no longer present at all in current OpenAI models, and capitulation manifests instead as a silent confident rewrite.

\para{Position.} We sit at the intersection of these three threads. We import the ladder design from the framing literature, the \flipflop two-turn protocol from the sycophancy literature, and the tone-as-stimulus question from the prompt-engineering literature, and combine them in the first single experiment that varies preemptive tone over a graded ladder with paired token, accuracy, apology, and post-rebuttal-flip measurements.
